% !TEX program = lualatex
% ============================================================
% main.tex — selve samlingen
%
% Denne fila er struktur og dikt. Alle typografiske valg bor
% i diktstil.sty, som også laster babel og setter
% papirformatet.
%
% Kompiler med LuaLaTeX, ikke pdflatex:
%   latexmk -lualatex main
% ============================================================
\documentclass[12pt,twoside]{book}


% ============================================================
% LAST INN DIKTSTIL
% ============================================================

% hefte står på i alle variantene: samlingen fylles opp med
% blanke sider til et sidetall delelig med fire, så et falset
% hefte får bakside. Ta opsjonen bort for løse ark.

% ------------------ Trykk, 140 x 216 mm -----------
\usepackage[hefte]{diktstil}

% ------------------ Egen utskrift, A5 -------------
% Samme satsflate, så linje- og sidebrudd er som over.
% \usepackage[a5,hefte]{diktstil}

% ------------------ Sentrerte dikt ----------------
% Hvert dikt midtstilt etter sin lengste linje.
% \usepackage[sentrert,hefte]{diktstil}

% ------------------ Sidekorrektur -----------------
% Siden uendret, sentrert på A4 med skjæremerker.
% \usepackage[korrektur,hefte]{diktstil}


% ============================================================
% DOKUMENT
% ============================================================
\begin{document}

\title{Samlingens tittel}
\subtitle{Dikt}
\author{Forfatternavn}
\date{\the\year}


% ============================================================
% FORSTOFF
% ============================================================

% ------------------ Smusstittel -------------------
% Valgfri: bare tittelen på en egen side foran tittelsiden.
% \halftitlepage

% ------------------ Tittelside --------------------
\booktitlepage

% ------------------ Kolofon -----------------------
\colophonpage{%
  © \the\year\ Forfatternavn\\
  Satt med Garamond Libre i \LaTeX\\
  ISBN 978-82-000-0000-0}

% ------------------ Dedikasjon --------------------
\dedicationpage{Til noen.}

% ------------------ Innhold -----------------------
\poemtoc


% ============================================================
% DIKTENE
%
% Ett dikt-miljø per dikt. Hvert dikt begynner på en ny side.
% Bytt ut plassholderne under med dine egne.
% ============================================================

\del{I}

\begin{dikt}[Et dikt med tittel]
Tittelen står i klammene
etter begin.

Hver linje i kilden
er en linje på siden,
og en blank linje
gir en ny strofe.
\end{dikt}

\begin{dikt}
Et dikt uten tittel
begynner der tittelen ellers ville stått

og står i innholdet
med sin første linje.
\end{dikt}

\begin{dikt}[Innrykk]
Mellomrom først på linja
    rykker linja inn,
        et halvt gefirt for hvert,

og mellomrom   inne i linja   blir stående.
\end{dikt}

\begin{dikt}[En lang linje]
En linje kan være kort.

En linje som er lengre enn satsflata er bred, brytes av seg selv, og fortsettelsen rykkes inn så den ikke kan forveksles med en ny linje.

Neste linje begynner ved venstrekanten igjen.
\end{dikt}


\del{II}

\begin{dikt}[Et langt dikt]
Dette diktet er lengre enn siden.
Det deles derfor i to,
men aldri midt i en strofe:
bruddet legges i lufta mellom to av dem.

Strofa som ikke får plass,
flyttes hel til neste side,
og da kan bunnen av siden
bli stående åpen.

Slik ser tredje strofe ut.
Den har fire linjer som de andre,
og sier ikke stort,
men den tar den plassen den skal.

Fjerde strofe følger etter.
Linjene står med seksten punkts avstand,
og mellom strofene
er det én linjes luft.

Femte strofe er den siste på denne siden.
Satsflata rommer tjueni linjer,
og diktet begynner på den sjette,
med tittelen på den fjerde.

Sjette strofe står øverst på neste side.
Fortsettelsen begynner øverst i satsflata,
uten lufta som står over tittelen,
siden det ikke er et nytt dikt.

Sidetallet står midtstilt nederst,
ett steg mindre enn teksten,
og det er ingen topptekst
over diktene.

Åttende strofe er den siste.
Neste dikt begynner på en ny side,
i samme høyde som dette begynte,
og slik fortsetter samlingen.
\end{dikt}

\begin{dikt}
Bytt ut disse diktene med dine egne.

Det er alt som skal til.
\end{dikt}


\end{document}
